\chapter{Modular Powers}

\section{The Multiplicative Order of an Integer Modulo $n$}

This section is about the function $n\mapsto a^n-1$. We will begin by showing that this function preserves divisibility.

\begin{proposition}
  Suppose that $m\mid n$. Then we have
  \begin{equation*}
    a^m-1\mid a^n-1.
  \end{equation*}
\end{proposition}

\begin{proof}
  Suppose that $qm=n$. Then we can write
  \begin{equation*}
    a^n-1 = (a^m)^{q} -1^q = (a^m-1)\sum_{k=0}^{q-1}a^{km}
  \end{equation*}
  by the formula for the difference of $q$th powers (\cref{exercise:difference-of-powers}).
\end{proof}

We will now define the multiplicative order of an element $a$ modulo $n$. If $\gcd(a,n)=1$, then we can simply define the multiplicative order of $a$ to be the least positive integer $k$ such that
\begin{equation*}
  a^k\equiv 1\pmod{n}.
\end{equation*}
However, if $\gcd(a,n)\neq 1$ then there will not be such a positive integer. In this case we set the multiplicative order of $a$ modulo $n$ to be $0$. At first glance, this may look quite arbitrary. However, there is a natural way of looking at it.

\begin{definition}
  Consider an integer $a$ and a natural number $n$. We define he \emph{multiplicative order} of $a$ modulo $n$ to be the unique natural number $m$ such that the ideal
  \begin{equation*}
    I_{a,n}:=\{k\in\mathbb{Z}\mid a^{|k|} \equiv 1\pmod{n}\}
  \end{equation*}
  is the ideal $(m)$ generated by $m$.
\end{definition}

Using this definition, we see that if $\gcd(a,n)=1$, then indeed the order $\ord_n(a)$ of $a$ modulo $n$ is the least positive integer $k$ such that $a^k\equiv 1\pmod{n}$. If $\gcd(a,n)\neq 1$, then the ideal $I_{a,n}$ is the zero ideal $(0)$, which implies that $\ord_n(a)=0$.

We saw the multiplicative order featuring in the second proof of \crefnamed{theorem:eulers-theorem}{Euler's Theorem}. Thus, using the second proof of Euler's theorem we immediately observe that
\begin{equation*}
  \ord_n(a)\mid \phi(n),
\end{equation*}
for any $a$ and $n$ such that $\gcd(a,n)=1$.

In the following theorem, which is also known as the \emph{Order Theorem}, we will see that by setting the value of the order of $a$ modulo $n$ to be $0$ when $\gcd(a,n)\neq 1$, we prove a powerful duality principle between multiples of $\ord_n(a)$ and numbers dividing $a^k-1$.

\begin{theorem}[Order Theorem]\label{theorem:order-theorem}
  The multiplicative order of an integer $a$ modulo $n$ satisfies the logical equivalence
  \begin{equation*}
    \ord_n(a)\mid k \Leftrightarrow n \mid a^k-1
  \end{equation*}
\end{theorem}

\begin{proof}
  We have two cases to consider: either $\gcd(a,n)=1$ or $\gcd(a,n)\neq 1$. In the first case, $\ord_n(a)$ is the least positive integer $m$ such that
  \begin{equation*}
    a^m\equiv 1 \pmod{n}.
  \end{equation*}
  It is immediate from this condition that if $m\mid k$, then it follows that $a^k\equiv 1\pmod{n}$. Conversely, if $a^k\equiv 1\pmod{n}$, then we can write $k=qm+r$ by the Euclidean division theorem, where $0\leq r<m$. It follows that
  \begin{equation*}
    a^r \equiv a^{qm}a^r \equiv a^{qm+r} \equiv a^k \equiv 1\pmod{n}. 
  \end{equation*}
  Since $r$ is strictly less than $m$ and $m$ is the least positive integer such that $a^m\equiv 1\pmod{n}$, it follows that $r=0$. Thus we conclude that $\ord_n(a)\mid k$. This proves the claim in the first case.

  In the second case we have $\ord_n(a)=0$. Then the condition that $\ord_n(a)\mid k$ is equivalent to the condition that $k=0$. Thus, our task is to show that $n\mid a^k-1$ if and only if $k=0$. Consider a common divisor $d>1$ of $a$ and $n$. Then we have $d\mid a^k-1$, so it follows that $d\mid a^k$ if and only if $d\mid 1$. Since we have assumed that $d>1$, it therefore follows that $d\nmid a^k$. However, since $d$ divides $a$, the only way in which this is possible is that $k=0$.
  \begin{equation*}
    \ord_n(a)\mid k.\qedhere
  \end{equation*}
\end{proof}

The previous theorem has strong implications about the function $k\mapsto a^k-1$, which we will present as a corollary of the following definition and theorem.

\begin{definition}
  A function $f:\mathbb{N}\to\mathbb{N}$ is said to be a \emph{divisibility sequence} if it preserves divisibility, i.e., if it satisfies the condition
  \begin{equation*}
    m\mid n \Rightarrow f(m)\mid f(n)
  \end{equation*}
  for all $m$ and $n$. A divisibility sequence is said to be a \emph{strong divisibility sequence} if $f(0)=0$ and $f$ preserves greatest common divisors, i.e., if
  \begin{equation*}
    f(\gcd(m,n))=\gcd(f(m),f(n))
  \end{equation*}
  for all $m$ and $n$.
\end{definition}

\begin{theorem}\label{theorem:galois-connections-divisibility}
  Suppose that $f:\mathbb{N}\to\mathbb{N}$ is a divisibility sequence. Then the following are equivalent:
  \begin{enumerate}
  \item There is a function $g:\mathbb{N}\to\mathbb{N}$ satisfying the logical equivalence
  \begin{equation*}
    g(m)\mid n \Leftrightarrow m\mid f(n).
  \end{equation*}
  \item The divisibility sequence $f$ is a strong divisibility sequence.
  \end{enumerate}
\end{theorem}

\begin{proof}
  Suppose first that a function $g$ exists satisfying the logical equivalence
  \begin{equation*}
    g(m)\mid n \Leftrightarrow m\mid f(n).
  \end{equation*}
  Then it follows that
  \begin{align*}
    d\mid f(\gcd(m,n)) & \Leftrightarrow g(d)\mid \gcd(m,n) \\
                       & \Leftrightarrow g(d)\mid m \land g(d)\mid n \\
                       & \Leftrightarrow d\mid f(m) \land d\mid f(n) \\
                       & \Leftrightarrow d\mid \gcd(f(m),f(n)).
  \end{align*}
  Since $f(\gcd(m,n))$ and $\gcd(f(m),f(n))$ have the exact same set of divisors, they must be equal.
  
  For the converse, suppose $f$ preserves greatest common divisors. Then consider the ideal
  \begin{equation*}
    I_m := \left\{ k_1n_1+\cdots+ k_rn_r \mathbin{\big|} m\mid f(n_i) \text{ for all } i \right\}
  \end{equation*}
  and define the function $g$ by letting $g(m)$ be the unique natural number such that $I_m=(g(m))$. In other words, $g(m)$ is the least natural number that can be written as a linear combination of numbers $n_i$, such that $m\mid f(n_i)$.

  To see that $g$ is a divisibility sequence, suppose that $m\mid m'$, i.e., suppose that $qm=m'$. Then any linear combination of natural numbers $n$ such that $m'\mid f(n)$ is also a linear combination of natural numbers $n$ such that $m\mid f(n)$. In particular $g(m')$ is in the ideal $(g(m))$, which means that $g(m)\mid g(m')$.

  We claim that the function $g$ satisfies the desired logical equivalence. Suppose first that $g(m)\mid n$, say $qg(m)=n$. Note that $g(m)$ can be written in the form $k_1n_1+\cdots+ k_rn_r$, so we can proceed by induction on $r$. In the base case, we have that $g(m)=0$. Then it follows that $n=0$, and it also follows that every linear combination of natural numbers $n$ such that $m\mid f(n)$ must be $0$. Since $f(0)=0$ by the assumption that $f$ is a strong divisibility sequence, this shows that $m\mid f(n)$.

  For the inductive step, suppose that for every $m$ such that $g(m)$ can be written in the form $k_1n_1+\cdots+ k_rn_r$, we have $g(m)\mid n$ implies $m\mid f(n)$, and consider $m$ such that $g(m)$ can be written in the form
  \begin{equation*}
    k_1n_1+\cdots+ k_rn_r + k_{r+1}n_{r+1}.\qedhere
  \end{equation*}
\end{proof}

\begin{corollary}
  Consider an integer $a$ and two natural numbers $m$ and $n$ relatively prime to $a$. Then
  \begin{equation*}
    \gcd(a^m-1,a^n-1)=a^{\gcd(m,n)}-1.
  \end{equation*}
\end{corollary}

The logical equivalence in the definition of the order of an element can be used effectively to compute the order of a power of an element in terms of the order of that element.

\begin{proposition}\label{proposition:order-power}
  For any integer $a$ of order $k=\ord_n(a)$, we have
  \begin{equation*}
    \ord_n(a^m)=\frac{k}{\gcd(m,k)}
  \end{equation*}
  In particular, we have $\ord_n(a^m)=k$ if and only if $m$ and $k$ are relatively prime.
\end{proposition}

\begin{proof}
  We have the following chain of logical equivalences:
  \begin{equation*}
    \ord_n(a^m) \mid l ~~\Leftrightarrow~~ n \mid (a^m)^l-1
    ~~\Leftrightarrow~~ \ord_n(a) \mid ml
    ~~\Leftrightarrow~~ \frac{\ord_n(a)}{\gcd(m,\ord_n(a))} \mid l.
  \end{equation*}
  Thus, we see that $\ord_n(a^m)$ and $k/\gcd(m,k)$ divide the same numbers, so they must be equal.
\end{proof}

\begin{remark}
  More generally, if $(P,\leq_P)$ and $(Q,\leq_Q)$ are two posets, a \emph{Galois connection} between $P$ and $Q$ consists of order preserving maps $f:P\to Q$ and $g:Q\to P$ satisfying the logical equivalence
  \begin{equation*}
    g(x)\leq_P y \Leftrightarrow x\leq_Q f(x)
  \end{equation*}
  for any $x\in Q$ and $y\in P$. In this case, the map $f$ is said to be the \emph{upper adjoint}, and the map $g$ is said to be the \emph{lower adjoint} of the Galois connection. Thus we have proven in this section that the function $n\mapsto a^n-1$ is the upper adjoint of a Galois connection, and we have used this fact to show that $f$ preserves greatest common divisors. \cref{theorem:galois-connections-divisibility} generalizes to the setting of arbitrary posets, and indeed, some readers may recognize it as a special case of the \emph{Adjoint Functor Theorem}.

  Another famous Galois connection on the natural numbers ordered by divisibility is the Fibonacci sequence. We will give the full statement of this fact in the exercises.
\end{remark}

\section{The Infinitude of Primes Congruent to 1 Modulo Powers of 2}

One way to show that a prime $p$ is congruent to $1$ modulo $n$ is by finding an integer $a$ such that
\begin{equation*}
  \ord_p(a)=n.
\end{equation*}
By the \crefnamed{theorem:order-theorem}{Order Theorem} and \crefnamed{theorem:fermats-little-theorem}{Fermat's Little Theorem}, this implies that $n\mid p-1$. We will use this idea to show that there are infinitely many primes congruent to $1$ modulo any power of $2$. In particular, there are infinitely many primes congruent to $1$ modulo $4$, modulo $8$, modulo $16$, and so on.

\begin{theorem}\label{theorem:infinitude-of-primes-1-mod-power-of-2}
  For any positive natural number $n$, there are infinitely many primes congruent to $1$ modulo $2^n$.
\end{theorem}

\begin{proof}
  Since $p\equiv 1\pmod{2^{n+1}}$ implies that $p\equiv 1\pmod{2^n}$, it suffices to show that for any finite list $p_1,\ldots,p_k$ of primes congruent to $1$ modulo $2^{n+1}$, there is a prime $q\equiv 1\pmod{2^{n+1}}$ which is not among the primes $p_i$.

  Given such a list of primes $p_i$, let
  \begin{equation*}
    a:=2p_1\cdots p_k,
  \end{equation*}
  and let $q$ be a prime divisor of the integer $b:=a^{2^n}+1$. Note that $q$ is necessarily an odd prime, since $a$ is even and hence $b$ is odd. Since $p_i\nmid b$ for any $i$, it follows that $q$ is not among the primes $p_i$. Furthermore, we observe that
  \begin{equation*}
    a^{2^n}\equiv -1\pmod{q}.
  \end{equation*}
  This implies that $\ord_q(a)\mid 2^{n+1}$, while clearly $\ord_q(a)\nmid 2^n$. Thus, we must have $\ord_q(a)=2^{n+1}$, and this allows us to conclude that $q\equiv 1\pmod{2^{n+1}}$.
\end{proof}

The previous theorem is a special case of a deep theorem in analytic number theory: \emph{Dirichlet's Theorem}. Dirichlet's Theorem shows that for any integer $a$ relatively prime to $n$, there are infinitely many primes
\begin{equation*}
  p\equiv a\pmod{n}.
\end{equation*}
The techniques of proving Dirichlet's Theorem are beyond the scope of this course. However, having taken this course, you are well equipped to read LeVeque's excellent \emph{Topics in Number Theory} \cite{leveque1956vol1,leveque1956vol2}, a two-volume work that contains well-explained proofs of Dirichlet's Theorem as well as the Prime Number Theorem.

%%% Local Variables:
%%% mode: latex
%%% TeX-master: "elementary-number-theory"
%%% End:
